\documentclass[12pt]{article}
\usepackage[a4paper, margin=1in]{geometry}
\usepackage{longtable}
\usepackage{titlesec}
\usepackage[table, svgnames, dvipsnames]{xcolor}
\usepackage{makecell, cellspace, caption}
\titleformat{\section}{\normalfont\Large\bfseries}{}{0em}{}

\title{\textbf{Course Schedule: Mechanics and Relativity WBPH001-10 2026:1a}}
\author{Based on David Morin's \textit{Introduction to Classical Mechanics}}
\date{}

\begin{document}

\maketitle

\section*{Overview}
\begin{itemize}
  \item \textbf{Lectures:} Mondays 09:00--11:00, Wednesdays 11:00--13:00
  \item \textbf{First Lecture:} Wednesday, September 2, 2026
  \item \textbf{Midterm Exam:} Thursday, October 1, 18:30--20:30 (Lectures 1--8)
  \item \textbf{Midterm Exam Review:} Date and time TBD (normally about one week after the exam)
  \item \textbf{Final Exam:} Monday, October 26, 18:15--20:15 (Comprehensive)
  \item \textbf{Final Exam Review:} Date and time TBD (normally about one week after the exam)
  \item \textbf{Resit Exam:} Date and time TBD (expected in February 2027)
  \item \textbf{Resit Exam Review:} Date and time TBD (normally about one week after the resit)
  \item \textbf{Textbook:} David Morin, \textit{Introduction to Classical Mechanics}
\end{itemize}


\newpage

\section*{Lecture Schedule}

\begin{longtable}{|c|c|c|p{5.8cm}|p{2.7cm}|}
\hline
\textbf{Lec.} & \textbf{Date} & \textbf{Time} & \textbf{Topic} & \textbf{Reading} \\ \hline
1 & Wed Sep 2 & 11--13 & Course overview; introduction and historical context & Intro / supplementary \\ \hline
2 & Mon Sep 7 & 09--11 & Foundations of special relativity: events, frames, and postulates & Ch. 11.1--11.3 \\ \hline
3 & Wed Sep 9 & 11--13 & Lorentz transformations and basic consequences & Ch. 11.4--11.5 \\ \hline
4 & Mon Sep 14 & 09--11 & The invariant interval, proper time, spacetime diagrams & Ch. 11.6--11.7 \\ \hline
5 & Wed Sep 16 & 11--13 & The Doppler effect, rapidity, and the pole-barn paradox & Ch. 11.8--11.10 \\ \hline
6 & Mon Sep 21 & 09--11 & Chapter 11 recap: spacetime diagrams and two-observer problems & Ch. 11 \\ \hline
7 & Wed Sep 23 & 11--13 & Dimensional analysis and natural units & Supplementary \\ \hline
8 & Mon Sep 28 & 09--11 & Fermi problems and estimation techniques & Supplementary \\ \hline
9 & Wed Sep 30 & 11--13 & Newtonian dynamics review: forces, energy, and momentum & Ch. 2, 3, 5 \\ \hline
-- & Thu Oct 1 & 18:30--20:30 & \multicolumn{2}{p{8.5cm}|}{\textbf{Midterm Exam (Lectures 1--8)}} \\ \hline
-- & TBD & TBD & \multicolumn{2}{p{8.5cm}|}{\textbf{Midterm Exam Review} (normally about one week after the exam)} \\ \hline
10 & Mon Oct 5 & 09--11 & Relativistic energy, momentum, and transformation rules &  Ch. 12.1--12.2 \\ \hline
11 & Wed Oct 7 & 11--13 & Collisions, decays, and particle-physics units & Ch. 12.3--12.4 \\ \hline
12 & Mon Oct 12 & 09--11 & Relativistic force and rocket motion & Ch. 12.5--12.6 \\ \hline
13 & Wed Oct 14 & 11--13 & 4-vectors: definitions, examples, and properties & Ch. 13.1--13.3 \\ \hline
14 & Mon Oct 19 & 09--11 & 4-vectors continued: energy, momentum, force, and physical laws & Ch. 13.4--13.6 \\ \hline
15 & Wed Oct 21 & 11--13 & Final comments: equivalence principle and accelerated frames & Ch. 14 \\ \hline
-- & Mon Oct 26 & 18:15--20:15 & \multicolumn{2}{p{8.5cm}|}{\textbf{Final Exam (Comprehensive)}} \\ \hline
-- & TBD & TBD & \multicolumn{2}{p{8.5cm}|}{\textbf{Final Exam Review} (normally about one week after the exam)} \\ \hline
-- & Feb 2027 & TBD & \multicolumn{2}{p{8.5cm}|}{\textbf{Resit Exam} (exact date TBD)} \\ \hline
-- & TBD & TBD & \multicolumn{2}{p{8.5cm}|}{\textbf{Resit Exam Review} (normally about one week after the resit)} \\ \hline
\end{longtable}

\end{document}
